\noindent\textbf{Part 1 (13 points): Cosmological Surface Brightness Dimming}
\begin{description}
\item[(a)] Consider a photon emitted with energy $E_{\mathrm{em}}$. Due to
  cosmic expansion, the energy received by the observer is redshifted,
  \[ E_{\mathrm{obs}} = \frac{E_{\mathrm{em}}}{1+z} \]
  \hfill(2 points)

  Due to cosmological time dilation, an emission time interval
  $\mathrm{d}t_{\mathrm{em}}$ is observed as $\mathrm{d}t_{\mathrm{obs}}$ at
  the receiver,
  \[ \mathrm{d}t_{\mathrm{obs}} = (1+z)\,\mathrm{d}t_{\mathrm{em}} \]
  \hfill(2 points)

  Combining these two effects, the observed luminosity is attenuated as
  \[ L_{\mathrm{obs}} = \frac{L}{(1+z)^2} \]
  By expressing the observed flux $F_{\mathrm{obs}}$ via the comoving distance
  $D_M$
  \[ F_{\mathrm{obs}} = \frac{L_{\mathrm{obs}}}{4\pi D_M^2}
     = \frac{L}{4\pi D_M^2 (1+z)^2} \]
  \hfill(1 point)

  and then comparing this to the definition
  \[ F_{\mathrm{obs}} = \frac{L}{4\pi D_L^2} \]
  we have
  \[ D_L = D_M (1+z) \]
  \hfill(1 point)

  The angular size of the object is
  \[ \theta = \frac{R(1+z)}{D_M} \]
  \hfill(2 points)

  By definition,
  \[ \theta = \frac{R}{D_A} \]
  Therefore,
  \[ D_A = \frac{D_M}{1+z} \]
  \hfill(1 point)

  Substituting $D_M = D_A(1+z)$ into the $D_L$ expression yields
  \[ D_L = D_A (1+z)^2 \]
  \hfill(1 point)

\item[(b)] To establish the equation for $\Sigma_{\mathrm{obs}}$, we start
  from the definition
  \[ \Sigma_{\mathrm{obs}} = \frac{F_{\mathrm{obs}}}{\Omega_{\mathrm{obs}}}
     = \frac{\dfrac{L}{4\pi D_L^2}}{\dfrac{\pi R^2}{D_A^2}} \]
  \hfill(1 point)

  By re-arranging the terms, we obtain
  \[ \Sigma_{\mathrm{obs}} = \frac{L}{\pi R^2} \cdot \frac{D_A^2}{4\pi D_L^2}
     = \Sigma_{\mathrm{int}} \cdot \frac{D_A^2}{4\pi D_L^2} \]
  \hfill(1 point)

  Substituting $D_A^2/D_L^2 = 1/(1+z)^4$ into the equation gives us
  \[ \Sigma_{\mathrm{obs}} = \frac{\Sigma_{\mathrm{int}}}{4\pi(1+z)^4} \]
  \hfill(1 point)
\end{description}

\medskip\noindent\textbf{Part 2 (8 points): Dissecting JWST Data}

The intrinsic rest-frame optical flux (dust-corrected) is given by
\[ F_{\mathrm{opt,int}} = F_{\mathrm{opt,obs}} \cdot 10^{\frac{A_V}{2.5}} \]
\hfill(1 point)

Substituting $F_{\mathrm{opt,obs}}$ and $A_V$ values into the equation
provides
\[ F_{\mathrm{opt,int}} = 5.72 \times 10^{-19} \times 10^{\frac{2.79}{2.5}}
   = 7.47 \times 10^{-18}\,\mathrm{W\,m^{-2}} \]
\hfill(1 point)

The intrinsic bolometric flux:
\[ F_{\mathrm{bol}} = F_{\mathrm{opt,int}} \cdot BC
   = 7.47 \times 10^{-18} \times 9.26
   = 6.92 \times 10^{-17}\,\mathrm{W\,m^{-2}} \]
\hfill(1 point)

In addition, after converting $\theta$ to radians
\[ \theta_{\mathrm{rad}} = \frac{0.05}{3600} \times \frac{\pi}{180}
   = 2.424 \times 10^{-7}\,\mathrm{rad} \]
\hfill(1 point)

we can calculate the observed solid angle by
\[ \Omega_{\mathrm{obs}} = \pi \theta_{\mathrm{rad}}^2
   = 1.846 \times 10^{-13}\,\mathrm{sr} \]
\hfill(1 point)

Therefore,
\[ \Sigma_{\mathrm{obs\_corr}} = \frac{F_{\mathrm{bol}}}{\Omega_{\mathrm{obs}}}
   = \frac{6.92 \times 10^{-17}}{1.846 \times 10^{-13}}
   = 3.75 \times 10^{-4}\,\mathrm{W\,m^{-2}\,sr^{-1}} \]
\hfill(1 point)

Applying the result from Part 1 gives us $\Sigma_{\mathrm{int}}$ in standard
SI units as
\[ \Sigma_{\mathrm{int\_SI}} = 4\pi(1+z)^4 \Sigma_{\mathrm{obs\_corr}} \]
Since $z = 7.26$, we have $(1+z)^4 \approx 4655.0$. Substituting this value
into the equation yields
\[ \Sigma_{\mathrm{int\_SI}} = 4 \times \pi \times 4655.0 \times 3.75 \times 10^{-4}
   = 21.94\,\mathrm{W\,m^{-2}} \]
\hfill(1 point)

We use the following conversion factor
\[ 1\,\mathrm{kpc^2} = (3.086 \times 10^{19}\,\mathrm{m})^2
   = 9.523 \times 10^{38}\,\mathrm{m^2} \]
Hence,
\[ \Sigma_{\mathrm{int}} = \frac{21.94 \times 9.523 \times 10^{38}}{3.828 \times 10^{26}}
   = 5.46 \times 10^{13}\,\mathrm{L_\odot\,kpc^{-2}} \]
\hfill(1 point)

\begin{center}
\begin{tabular}{ccc}
\toprule
Half credit lower limit & Full credit range & Half credit upper limit \\
\midrule
$5.35 \times 10^{13}\,\mathrm{L_\odot\,kpc^{-2}}$ &
  \begin{tabular}[t]{@{}c@{}}$5.40 \times 10^{13}$ to\\
  $5.50 \times 10^{13}\,\mathrm{L_\odot\,kpc^{-2}}$\end{tabular} &
  $5.55 \times 10^{13}\,\mathrm{L_\odot\,kpc^{-2}}$ \\
\bottomrule
\end{tabular}
\end{center}

\medskip\noindent\textbf{Part 3 (18 points): Eddington Dynamical Model for the Star-Forming Disk}
\begin{description}
\item[(a)] The Eddington surface luminosity density is
  \[ \Sigma_{\mathrm{int\_max}} = \frac{L_{\mathrm{Edd}}}{\pi R^2} \]
  The total radiation force exerted on a dust particle is
  \[ f = \sigma \frac{\Sigma_{\mathrm{int\_max}}}{c} \]
  \hfill(2 points)

  where $\sigma$ is the dust cross-section given by
  \[ \sigma = m_{\mathrm{d}} \kappa_{\mathrm{dust}} \]
  \hfill(1 point)

  where $m_{\mathrm{d}}$ is the mass of dust particle.

  The acceleration driven by radiation pressure is
  \[ g_{\mathrm{rad}} = \frac{f}{m_{\mathrm{d}}}
     = \frac{\kappa_{\mathrm{dust}} \Sigma_{\mathrm{int\_max}}}{c} \]
  \hfill(1 point)

  Based on the Gauss's gravity theorem, we have
  \[ 2\pi R^2 g_z = 4\pi G M_{\mathrm{tot}} \]
  \hfill(1 point)

  where $M_{\mathrm{tot}}$ is the total disk's mass.

  Hence,
  \[ g_z = 2\pi G \Sigma_{\mathrm{tot}} \]
  \hfill(1 point)

\item[(b)] Based on the balance between radiation pressure and gravity, we
  can establish the following critical equilibrium condition
  \[ g_{\mathrm{rad}} = g_z \]
  \hfill(1 point)
  \begin{align*}
    \Rightarrow \frac{\kappa_{\mathrm{dust}} \Sigma_{\mathrm{int\_max}}}{c}
      &= 2\pi G \Sigma_{\mathrm{tot}} \\
    \Rightarrow \Sigma_{\mathrm{int\_max}}
      &= \frac{2\pi G c \Sigma_{\mathrm{tot}}}{\kappa_{\mathrm{dust}}}
  \end{align*}
  \hfill(1 point)

  The Eddington SFR surface density is derived from the Kennicutt scaling
  law:
  \[ \Sigma_{\mathrm{Edd}} = \frac{\Sigma_{\mathrm{int\_max}}}{k} \]
  \hfill(1 point)
  \[ \Rightarrow \Sigma_{\mathrm{Edd}}
     = \frac{2\pi G c \Sigma_{\mathrm{tot}}}{k \kappa_{\mathrm{dust}}} \]
  \hfill(1 point)

  Converting $\Sigma_{\mathrm{tot}}$ to SI units gives us
  \[ \Sigma_{\mathrm{tot}}
     = \frac{2.3 \times 10^{10} \times 1.988 \times 10^{30}}{9.523 \times 10^{38}}
     = 48.01\,\mathrm{kg\,m^{-2}} \]
  \hfill(0.5 points)

  In SI units, $\Sigma_{\mathrm{int\_max}}$ is calculated by
  \[ \Sigma_{\mathrm{int\_max}}
     = \frac{2\pi \times 6.674 \times 10^{-11} \times 2.998 \times 10^{8} \times 48.01}{1.5}
     = 4.02\,\mathrm{W\,m^{-2}} \]
  \hfill(1 point)

  Using unit conversion
  \[ 4.02\,\mathrm{W\,m^{-2}} = 1.00 \times 10^{13}\,\mathrm{L_\odot\,kpc^{-2}} \]
  \hfill(0.5 points)

  we obtain
  \[ \Sigma_{\mathrm{Edd}} = \frac{1.00 \times 10^{13}}{10^{10}}
     = 1000\,\mathrm{M_\odot\,yr^{-1}\,kpc^{-2}} \]
  \hfill(1 point)

  \begin{center}
  \begin{tabular}{ccc}
  \toprule
  Half credit lower limit & Full credit range & Half credit upper limit \\
  \midrule
  980 & 990 to 1010 & 1020 \\
  \bottomrule
  \end{tabular}
  \end{center}

\item[(c)] Using $\Sigma_{\mathrm{int}}$ in Part 2 provides
  \[ \Sigma_{\mathrm{SFR\_actual}} = \frac{\Sigma_{\mathrm{int}}}{k} \]
  \hfill(1 point)
  \[ \Sigma_{\mathrm{SFR\_actual}}
     = \frac{5.46 \times 10^{13}}{1.0 \times 10^{10}}
     = 5460\,\mathrm{M_\odot\,yr^{-1}\,kpc^{-2}} \]
  \hfill(1 point)

  We can compare $\Sigma_{\mathrm{SFR\_actual}}$ with $\Sigma_{\mathrm{Edd}}$,
  \[ \frac{\Sigma_{\mathrm{SFR\_actual}}}{\Sigma_{\mathrm{Edd}}} = 5.46 \]
  \hfill(1 point)

  \emph{It is clear to see that the actual density vastly exceeds the
  Eddington limit. The disk is highly super-Eddington.}

  The system is dynamically unstable. \hfill(2 points)

  \emph{If the luminosity were purely stellar in origin, the extreme
  radiation pressure would immediately unbind the gas and dust, triggering a
  catastrophic blowout and destroying the disk. Thus, the assumption of a
  purely starburst-driven luminosity is unphysical/untenable.}
\end{description}

\medskip\noindent\textbf{Part 4 (11 points): Clues from High-Resolution Spectroscopy}
\begin{description}
\item[(a)] Applying the Virial theorem to the broad-line region gives us
  \[ M = \frac{v^2 R_{\mathrm{em}}}{G} \]
  \hfill(2 points)

  On that basis, we can convert units and compute
  \begin{align*}
    M &= \frac{(3.4 \times 10^6)^2 \times (0.18 \times 3.086 \times 10^{16})}{6.674 \times 10^{-11}}
       = 9.62 \times 10^{38}\,\mathrm{kg} \\
    \Rightarrow M &= 4.84 \times 10^{8}\,\mathrm{M}_\odot
  \end{align*}
  \hfill(1 point)

  The object is a \textbf{supermassive black hole (C)}. \hfill(1 point)

  \emph{A huge mass of nearly half a billion solar masses confined within a
  sub-parsec radius ($< 0.2$\,pc) is a definitive signature of a Supermassive
  Black Hole (SMBH) / Active Galactic Nucleus (AGN).}

\item[(b)] \emph{Spectral analysis reveals a ``highly ionized, intensely
  hot'' gas phase. Under this extreme thermal and radiation environment, the
  interstellar dust grains are completely sublimated (destroyed).}

  \emph{Without dust, the radiation pressure now couples solely to free
  electrons via Thomson scattering. Consequently, the effective opacity
  $\kappa$ precipitously drops from
  $\kappa_{dust} = 1.5\,\mathrm{m^2\,kg^{-1}}$ to
  $\kappa_e = 0.04\,\mathrm{m^2\,kg^{-1}}$.}

  Assume the same optical thick disk model as above
  \[ L_{\mathrm{Edd}} = \frac{2\pi G M c}{\kappa_{\mathrm{e}}} \]
  \hfill(2 points)
  \[ \Sigma_{\mathrm{Edd}} = \frac{2 G M c}{\kappa_{\mathrm{e}} R^2}
     \approx 37500\,\mathrm{M_\odot\,yr^{-1}\,kpc^{-2}} \]
  % sic: the printed formula omits the division by k, and 37500 does not
  % follow from it with R = R_em (that gives about 7.8e9 M_sun/yr/kpc^2);
  % 37500 is the Part 3b value 1000 scaled by kappa_dust/kappa_e = 37.5,
  % as the Note below states.
  \hfill(2 points)
  \[ \frac{\Sigma_{\mathrm{SFR\_actual}}}{\Sigma_{\mathrm{Edd}}}
     = \frac{5460}{37500} = 0.15 \]
  \hfill(1 point)

  \emph{\underline{Note:} The transition to Thomson opacity
  ($\kappa_{\mathrm{e}}$) increases the medium's threshold for radiation
  pressure by a factor of $1.5/0.04 = 37.5$. Hence,
  $\dfrac{\Sigma_{SFR\_actual}}{\Sigma_{Edd}} = \dfrac{5.46}{37.5} = 0.15$.
  Award full points.}

  Since $\Sigma_{\mathrm{SFR\_actual}} < \Sigma_{\mathrm{Edd}}$, the radiation
  pressure does not exceed the gravitational pull of the SMBH. The system
  effectively maintains hydrostatic equilibrium. This system is radiatively
  stable (YES). \hfill(2 points)
\end{description}
